% Answer keys, teacher rubrics and scoring for the Level 1 assessments
% (input by answer-key.tex; the student papers are in assessments.tex)
\kcpart{Assessments}{Assessment Answer Keys}{Survey guidance, checkpoint and exit answers, teacher rubrics and scoring.}

\begin{teachernote}[Principles for every assessment]
\begin{itemize}[itemsep=1pt]
  \item \textbf{Nothing here is a grade to report.} The results tell you what to re-teach and who needs the
        support or extension path. Share results with the student as \enquote{you can now\dots{}} and \enquote{next we will\dots}.
  \item \textbf{Read questions aloud on request}, accept spoken answers, and give extra time. We are
        checking coding ideas, not reading speed, handwriting or spelling.
  \item \textbf{Accept any correct answer}, even one that is not printed here. If a student explains a
        different answer that makes sense, give the point.
  \item Record the band (Just starting / Getting there / Got it / Expert) in each student's progress tracker.
\end{itemize}
\end{teachernote}

\newcommand{\bandtable}[4]{%
  \begin{center}\small\renewcommand{\arraystretch}{1.3}
  \begin{tabular}{@{}llll@{}}\toprule
  \faSeedling\ Just starting & \faShoePrints\ Getting there & \faCheckCircle\ Got it & \faRocket\ Expert\\\midrule
  #1 & #2 & #3 & #4\\\bottomrule
  \end{tabular}\end{center}}

%% ---------------------------------------------------------------------
\section{Week 1: Interest and Comfort Survey}
There are no right answers and no score. Spend five minutes after the session reading the surveys and use them like this:

\begin{center}\small\renewcommand{\arraystretch}{1.35}
\begin{tabularx}{\linewidth}{@{}>{\bfseries\RaggedRight}p{3.4cm}X@{}}
\toprule
Survey section & What to do with it\\
\midrule
1. What do I like? & Use it to choose examples and free-build suggestions: a student who loves animals gets the \enquote{fish} or \enquote{cat} version of a task; game fans get the score and catch-game twists in weeks 3--4.\\
2. How do I feel? & \enquote{Worried} or \enquote{Not sure}: pair them as navigator first with a calm partner, praise effort out loud in week 1, and check in privately at the end of week 2.\\
3. Where am I starting? & Rows ticked \enquote{Just starting} for mouse or typing: give extra drag-and-drop time in Scratch (week 1--2) and a typing warm-up before week 7. Keep this page: students compare it with the week-12 self-rating.\\
4. How do I like to learn? & \enquote{Bigger letters}: print the dyslexia-friendly edition and turn up Thonny's font size. \enquote{Read aloud}: read instructions aloud for the whole group. \enquote{Short breaks}: use the stretch break from week 2 every session. Use the driver/navigator preference for the first pairing.\\
5. My dream project & Mention it again in week 10 when they choose a project idea. Answer their question during the first weeks if you can.\\
\bottomrule
\end{tabularx}
\end{center}
Keep the surveys private (in your folder, not on the wall). They contain no personal data beyond a nickname.

%% ---------------------------------------------------------------------
\section{Week 6: Mid-Level Checkpoint (8 points)}

\subsection{Answers}
\begin{answer}[Part A: What will the sprite do? (4 points)]
\begin{description}[leftmargin=0pt,itemsep=3pt,font=\bfseries]
  \item[A1 (G2).] \textbf{A square.} Four times: move 100, turn 90 degrees. $4 \times 90 = 360$, a full turn.
  \item[A2 (G4).] \textbf{6.} The score starts at 0 and goes up by 2 three times: 2, 4, 6.
        A common wrong answer is \enquote{2} (thinking the loop runs once) or \enquote{3} (counting the repeats): ask them to trace it with a finger.
  \item[A3 (G3).] \textbf{When the space key is pressed.} Not when the green flag is clicked: that is a different event hat block.
  \item[A4 (G3).] \textbf{\enquote{OK, dogs then!}} The answer \texttt{no} is not equal to \texttt{yes}, so the \textbf{else} part runs.
\end{description}
\end{answer}

\begin{answer}[Part B: Put the steps in order (2 points)]
\textbf{B1 (G1).} 1 Fill the pot with soil. 2 Make a small hole in the soil with your finger.
3 Drop the seed into the hole. 4 Cover the seed with soil. 5 Water the pot. 6 Put the pot in a sunny place.\\
Accept 5 and 6 swapped (watering after moving the pot also works). Give the point for any order a robot could really follow:
the hole must come before the seed, and the seed before covering it.

\textbf{B2 (G2).} \enquote{\textbf{repeat 4: clap, stomp}, then \textbf{bow}}. The bow must be \textbf{outside} the repeat
(it happens once). Give the point for \enquote{repeat 4 (clap, stomp)} even if the bow is missing, but mention it.
\enquote{Repeat 8: clap, stomp} is not correct: ask them to act it out and count.
\end{answer}

\begin{answer}[Part C: Find the bug (2 points)]
\textbf{C1 (G6, G2).} The turn is wrong: \textbf{turn 90} should be \textbf{turn 120}. A triangle needs 3 turns that add up
to 360 degrees, and $360 \div 3 = 120$. (Changing \enquote{repeat 3} to \enquote{repeat 4} gives a square, not a triangle: no point.)

\textbf{C2 (G6, G5).} The message names do not match: the cat sends \textbf{start game} but the dog listens for
\textbf{begin}, so nobody receives \enquote{start game}. Fix: change the dog's hat block to \textbf{when I receive start game}
(or make the cat broadcast \textbf{begin}). Give the point for spotting the mismatch, even without the fix written.
\end{answer}

\needspace{10\baselineskip}\subsection{Scoring}
Add up the points (0--8) and record the band:
\bandtable{0--2 points}{3--4 points}{5--6 points}{7--8 points}
\textbf{After the checkpoint:} for every question most students missed, re-teach it in the week-7 warm-up
(5 minutes) or use buffer B1 (Unplugged Games Day) before starting Python. A student in \enquote{Just starting}
needs the support path in weeks 7--9 (see the Teacher Guide) and more navigator time with a patient partner.
A student in \enquote{Expert} gets the extension path and can be a \enquote{bug-hunt helper} in week 7.

\subsection{Teacher rubric: Scratch goals (weeks 1--6)}
Use the checkpoint \textbf{and} what you saw in sessions. One mark per row.
{\small\renewcommand{\arraystretch}{1.3}
\begin{longtable}{@{}>{\bfseries\RaggedRight}p{2.2cm}*{4}{>{\RaggedRight\arraybackslash}p{\dimexpr(\linewidth-2.2cm-8\tabcolsep)/4}}@{}}
\toprule
Goal & \faSeedling\ Just starting & \faShoePrints\ Getting there & \faCheckCircle\ Got it & \faRocket\ Expert\\
\midrule\endhead
\bottomrule\endfoot
G1 Algorithms & Steps often out of order; needs help to sequence. & Orders familiar steps; misses hidden steps. & Writes and orders steps a robot could follow (B1). & Spots ambiguous steps and makes them precise.\\\midrule
G2 Loops & Copies repeat blocks without predicting them. & Predicts simple repeats (A1). & Uses repeat to shorten code; knows what is inside vs after (B2). & Chooses repeat counts and angles with $360 \div n$ on their own.\\\midrule
G3 Events and conditionals & Uses only the green flag. & Uses key events; if-blocks with help. & Predicts if/else (A3, A4); uses touching/sensing. & Combines events and conditions in own games.\\\midrule
G4 Variables and random & Variable blocks used only when copied. & Sets/changes a score with help. & Traces a variable (A2); uses pick random. & Uses variables for new purposes (lives, timer, size).\\\midrule
G5 Decomposition & One long script. & Uses broadcast with help. & Uses broadcasts and a My Block; can name the parts of a program. & Makes My Blocks with inputs and reuses them.\\\midrule
G6 Debugging & Asks for help straight away. & Finds a bug when told where to look. & Finds and fixes bugs like C1, C2. & Explains a strategy (test small parts, check names, act it out).\\
\end{longtable}}

%% ---------------------------------------------------------------------
\section{Week 12: Exit Assessment (14 points)}

\subsection{Answers}
\begin{answer}[Part A: Predict (5 points)]
\begin{description}[leftmargin=0pt,itemsep=3pt,font=\bfseries]
  \item[A1 (G9, G2).] \textbf{A triangle} with three equal sides (100 steps). The turtle goes right along the bottom,
    then turns left, so the triangle sits \textbf{above} the starting point with its point at the top.
    Give the point for any closed triangle.
  \item[A2 (G9, G2).] Four lines: \texttt{Hip}, \texttt{Hip}, \texttt{Hip}, \texttt{Hooray!}. \texttt{Hooray!} is printed
    \textbf{once} because it is not indented, so it is not inside the loop. Writing \enquote{Hooray!} three times is
    the classic indentation misconception: re-teach with the \enquote{Indent the line-up} activity from week 8.
  \item[A3 (G3, G4), 2 points.] With \texttt{stars = 7}: \texttt{Wow, a starry night!} (1 point).
    With \texttt{stars = 5}: \texttt{Just a few stars.} (1 point), because \texttt{5 > 5} is False: 5 is not
    \emph{bigger} than 5.
  \item[A4 (G5).] \textbf{Two squares.} Both start at the same corner (the middle of the screen). The second one,
    \texttt{square(100)}, is bigger: its sides are twice as long, so the small square sits inside the corner of the big one.
\end{description}
\end{answer}

\begin{answer}[Part B: Fix the bugs (4 points)]
\textbf{B1 (G6), 1 point per bug:}
\begin{enumerate}[itemsep=1pt]
  \item Line 3: \texttt{input} gives text, so it needs \texttt{int}: \texttt{sides = int(input("How many sides? "))}.
        (Also correct: a new line \texttt{sides = int(sides)}.)
  \item Line 4: missing colon: \texttt{for i in range(sides):}
  \item Line 5: spelling: \texttt{t.forward(80)}
\end{enumerate}
The corrected program (tested: with the answer \texttt{6} it draws a closed hexagon):
\begin{pycode}
import turtle
t = turtle.Turtle()
sides = int(input("How many sides? "))
for i in range(sides):
    t.forward(80)
    t.left(360 / sides)
turtle.done()
\end{pycode}
If a student says \enquote{\texttt{360 / sides} is a bug too}, they are thinking well: it only breaks because
\texttt{sides} is text. Once bug 1 is fixed, it works. Give the point for bug 1 if they fix it there instead.

\textbf{B2 (G5, G6).} The function \texttt{flower} is \textbf{defined but never called}: a recipe nobody cooks.
Add a call \textbf{before} \texttt{turtle.done()} and \textbf{not indented}, for example
\texttt{flower("\#CC79A7")} (any colour works, e.g. \texttt{flower("purple")}).
\end{answer}

\begin{answer}[Part C: Build it (3 points)]
One point for each tick-box: \textbf{(1)} it asks the question with \texttt{input} and turns the answer into a number with
\texttt{int}; \textbf{(2)} a loop draws that many circles; \textbf{(3)} each circle is in a random place
(\texttt{penup}, \texttt{goto} with \texttt{random.randint}). Run the program with the answer \texttt{5} and count.
A sample answer (one of many):
\begin{pycode}[exit\_circles.py (sample answer)]
import random
import turtle

t = turtle.Turtle()
t.speed(0)

count = int(input("How many circles? "))
for i in range(count):
    x = random.randint(-300, 300)
    y = random.randint(-200, 200)
    t.penup()
    t.goto(x, y)
    t.pendown()
    t.circle(30)

turtle.done()
\end{pycode}
\textbf{Partial credit:} a program that draws a fixed number of circles at random places scores 2 (input missing);
circles all in one place score 2 (random missing). Lines between circles (no \texttt{penup}) do not lose a point but
are worth mentioning. The \enquote{extra} (own function, colours) earns no extra points but is evidence for
\enquote{Expert} in the rubric below.
\end{answer}

\begin{answer}[Part D: Explain it (2 points)]
\textbf{1 point:} \texttt{size} is the function's input (parameter): it decides how long each side is, so
\texttt{square(50)} is small and \texttt{square(100)} is big.
\textbf{1 point:} a function is useful because you write the steps once and use them many times (with different sizes),
so the program is shorter, easier to read, and a fix in one place fixes every square.
Accept kid language: \enquote{size is how big}, \enquote{you don't have to write the loop again}, \enquote{like a My Block}.
\end{answer}

\needspace{10\baselineskip}\subsection{Scoring}
Add up the points (0--14) and record the band, then combine it with the project rubric to decide on the certificate and next steps.
\bandtable{0--4 points}{5--8 points}{9--11 points}{12--14 points}
\begin{itemize}[itemsep=1pt]
  \item \textbf{Ready for Level 2:} \enquote{Got it} or \enquote{Expert} here, and a project that runs. Most students should be here.
  \item \textbf{Ready with support:} \enquote{Getting there}. Start Level 2 and use its week-1 review and support paths; revisit \texttt{for} loops and \texttt{input}/\texttt{int} in the first two weeks.
  \item \textbf{Needs more time:} \enquote{Just starting}. This is not a failure. Use B2 (Free-Build and Remix Studio) and repeat weeks 8--10 exercises before Level 2. Every student still gets the Level 1 certificate: it celebrates effort and the finished project.
\end{itemize}

\subsection{Teacher rubric: Python goals (weeks 7--12)}
{\small\renewcommand{\arraystretch}{1.3}
\begin{longtable}{@{}>{\bfseries\RaggedRight}p{2.2cm}*{4}{>{\RaggedRight\arraybackslash}p{\dimexpr(\linewidth-2.2cm-8\tabcolsep)/4}}@{}}
\toprule
Goal & \faSeedling\ Just starting & \faShoePrints\ Getting there & \faCheckCircle\ Got it & \faRocket\ Expert\\
\midrule\endhead
\bottomrule\endfoot
G9 Reading and writing Python & Types code only when dictated. & Reads short programs; predicts with help. & Predicts A1--A3 correctly; writes C1 with few errors. & Writes new programs from a blank file.\\\midrule
G2 Loops & Unsure what repeats. & Uses \texttt{for}/\texttt{range} but confuses indentation (A2). & Uses loops correctly, including indentation. & Uses nested loops or the loop variable creatively.\\\midrule
G3, G4 Input, variables, if & Uses \texttt{input} but forgets \texttt{int}. & Uses variables; \texttt{if} with help. & Uses \texttt{input}/\texttt{int}, \texttt{if/else}, comparisons (A3). & Uses \texttt{elif}, handles unexpected answers.\\\midrule
G5 Functions & Cannot yet say what \texttt{def} does. & Calls given functions. & Defines and calls a function with a parameter (A4, B2, D1). & Writes functions with several parameters and reuses them.\\\midrule
G6 Debugging & Waits for help. & Fixes one of the three bugs in B1. & Fixes all three; reads Thonny's error message. & Explains why each bug breaks the program.\\\midrule
G7 Explaining & Says what the program does. & Points to the parts. & Explains how a part works (D1). & Explains choices and a bug they fixed.\\
\end{longtable}}
